% --- Chapter: Introduction ---
\section{Introduction}

The landscape of scripting languages is dominated by a trade-off between expressiveness and safety. Languages like Perl offer powerful text processing and concise syntax but lack memory safety guarantees. Languages like Python provide safety through garbage collection but incur performance overhead. Rust offers safety and performance but has a steep learning curve and is not designed for scripting.

\textbf{CRISP} (Creative Rust Implemented Scripting Paradigm) addresses this gap. It provides:

\begin{itemize}
    \item \textbf{Perl-inspired syntax:} Regular expressions (\texttt{=~}, \texttt{s///}), sigils (\texttt{\$}, \texttt{@}, \texttt{\%}), and built-in functional tools (\texttt{map}, \texttt{filter}, \texttt{grep})
    \item \textbf{Rust-powered safety:} Memory safety, no segmentation faults, use-after-free protection
    \item \textbf{Scripting agility:} REPL, one-liners, quick prototyping, file execution
    \item \textbf{System access:} Full POSIX system calls via opt-in \texttt{use posix;} module
    \item \textbf{Modern error handling:} \texttt{try}/\texttt{catch}/\texttt{finally} with \texttt{throw}/\texttt{die} keywords
    \item \textbf{Pattern matching:} Rust-style \texttt{match} with \texttt{where} guards
    \item \textbf{Object-oriented programming:} Classes, constructors, inheritance, static methods, and \texttt{super} calls
\end{itemize}

\begin{tcolorbox}[colback=lightlavender, colframe=softvioletdark, title=\textbf{\textcolor{darkpurple}{Definition: CRISP}}, fonttitle=\bfseries, sharp corners]
\textbf{CRISP} is a tree-walking interpreted scripting language implemented in Rust, characterized by:
\begin{itemize}
    \item Perl-inspired syntax with optional sigils (\texttt{\$scalar}, \texttt{@array}, \texttt{\%hash}, \texttt{\&ref})
    \item Native regular expressions with match (\texttt{=~}) and substitution (\texttt{s///})
    \item Lambda functions with explicit \texttt{return} and expression bodies
    \item Automatic hash key conversion and dot-notation field access
    \item Full POSIX system call integration (opt-in, namespaced)
    \item Console I/O with type-safe input functions and proper flush semantics
    \item Loop control flow with \texttt{break} and \texttt{continue}
    \item Variable assignment without redeclaration (\texttt{x = x + 1})
    \item Classes with constructors, inheritance, and static methods
\end{itemize}
\end{tcolorbox}
